\documentclass[12pt,a4paper]{article}
\usepackage[brazilian]{babel}
\usepackage[utf8]{inputenc}
\usepackage{amsmath,amssymb,amsthm}
\usepackage{geometry}
\usepackage{enumitem}
\usepackage{hyperref}
\usepackage{xcolor}
\usepackage{graphicx}
\usepackage{fancyhdr}
\usepackage{tcolorbox}
\usepackage{totpages}
\usepackage{titlesec}
\usepackage{etoolbox}
\usepackage{multicol}
\usepackage{multicol}
\setlength{\columnseprule}{0.4pt} % Espessura da linha divisória
\renewcommand{\columnseprulecolor}{\color{borderlight}} % Cor da linha (opcional)

\geometry{top=3.5cm, bottom=2.5cm, left=2.0cm, right=2.0cm, headheight=2.6cm}

\definecolor{primary}{RGB}{10, 45, 90}
\definecolor{secondary}{RGB}{25, 110, 180}
\definecolor{bglight}{RGB}{245, 247, 250}
\definecolor{borderlight}{RGB}{210, 220, 230}
\definecolor{correctgreen}{RGB}{40, 140, 60}

\tcbuselibrary{skins, breakable}

\newtcolorbox{boxfixacao}[1]{
  enhanced,
  colback=bglight,
  colframe=primary,
  arc=3pt,
  boxrule=1pt,
  title={#1},
  fonttitle=\sffamily\bfseries\color{white},
  coltitle=white,
  colbacktitle=primary,
  breakable
}

\newcounter{numquestao}
\setcounter{numquestao}{0}

\newenvironment{questao}[2][]{%
  \refstepcounter{numquestao}%
  \vspace{0.3cm}%
  \begin{tcolorbox}[
    enhanced,
    colback=white,
    colframe=secondary,
    arc=2pt,
    boxrule=0.8pt,
    leftrule=4pt,
    title={\textbf{Questão \thenumquestao} \ifstrempty{#2}{}{--- #2} \hfill \small\textnormal{#1}},
    fonttitle=\sffamily\bfseries\color{white},
    coltitle=white,
    colbacktitle=secondary,
    breakable
  ]%
}{%
  \end{tcolorbox}%
}

\newenvironment{gabarito}{%
  \vspace{0.1cm}%
  \begin{tcolorbox}[
    enhanced,
    colback=bglight,
    colframe=correctgreen,
    arc=2pt,
    boxrule=0.6pt,
    title={\small\sffamily\bfseries Resolução / Gabarito},
    fonttitle=\sffamily\color{white},
    colbacktitle=correctgreen,
    breakable
  ]%
  \small
}{%
  \end{tcolorbox}%
}

\newcommand{\espacoresposta}[1]{%
  \vspace{#1}%
}

\pagestyle{fancy}
\fancyhf{}

\fancyhead[L]{%
  \small\sffamily
  \textbf{IFPE -- Campus Caruaru}
}
\fancyhead[R]{%
  \small\sffamily
  \textbf{Instituto Federal de Pernambuco -- Campus Caruaru} \\
  \textcolor{secondary}{Disciplina:} Cálculo 1 \\
  \textcolor{secondary}{Prof.:} Cleibson José da Silva
}

\fancyfoot[L]{\small\sffamily\color{gray}Lista 2 -- Limites Trigonométricos}
\fancyfoot[R]{\small\sffamily\color{gray}Página \thepage\ de \ref{TotPages}}

\renewcommand{\headrulewidth}{0.8pt}
\renewcommand{\headrule}{{\color{primary}\hrule width\headwidth height \headrulewidth}}
\renewcommand{\footrulewidth}{0.4pt}

\hypersetup{
    colorlinks=true,
    linkcolor=secondary,
    urlcolor=secondary,
    citecolor=primary
}

\begin{document}

\begin{center}
  {\LARGE\bfseries\color{primary} Lista 3: Definição de Derivada} \\[0.1cm]
  {\small\sffamily\color{gray} Exercícios sobre definição de Derivada}
\end{center}

\vspace{0.2cm}


\section*{Lista 03 - Definição de Derivada}



\textbf{Use a definição:} \( f'(x) = \lim\limits_{h \to 0} \dfrac{f(x+h) - f(x)}{h} \) para resolver os problemas de 1 a 4.\\[0.5cm]

\begin{enumerate}
    \item  Dada \( f(x) = 4x - 3 \), encontre \( f'(2) \) usando a definição.
    
    \item  Dada \( f(x) = x^2 + 1 \), encontre \( f'(1) \) usando a definição.
    
    \item  Dada \( f(x) = \dfrac{1}{x} \), encontre \( f'(2) \) usando a definição.
    
    \item  Dada \( f(x) = 3x + 5 \), use a definição para encontrar a função derivada \( f'(x) \).
    \item  Um objeto se move segundo a função posição \( s(t) = t^2 + 3t \) (\( s \) em metros, \( t \) em segundos).
    
    \begin{enumerate}
        \item Qual sua velocidade média entre \( t = 1 \) e \( t = 4 \)?
        \item Encontre a velocidade instantânea em \( t = 2 \) s usando a definição de derivada.
    \end{enumerate}

    \item Uma população de bactérias é dada por \( P(t) = 200 + 30t \) (\( t \) em horas).
    Calcule a taxa de crescimento instantânea para qualquer tempo \( t \) usando a definição de derivada.

    \item Para \( f(x) = x^3 - 2x \):
    \begin{enumerate}
        \item Encontre \( f'(x) \)
        \item Calcule \( f'(1) \) usando o resultado do item (a)
        \item Escreva a equação da reta tangente ao gráfico de \( f \) no ponto onde \( x = 1 \)
    \end{enumerate}

  \end{enumerate}

  
\end{document}